\documentclass[12pt,a4paper]{article}

\usepackage{amsmath, amssymb, graphicx}
\usepackage{geometry}
\usepackage{setspace}
\usepackage{hyperref}

\geometry{margin=1in}
\setstretch{1.2}

\title{\textbf{Lecture Scribe}\\Fundamentals of Probability in Computing}
\author{
Name: Sakshi Patel\\
Student ID: AU2440206\\
Course: Fundamentals of Probability in Computing\\
}
\begin{document}
\date{}
\maketitle
\newpage

\section*{Step 1: Main Topics Identified (from the PDF)}

From the outline in the lecture (pages 2–3):

\begin{enumerate}
\item General Course Info + Why learn CSE400?

Engineering motivation

Applications: Speech recognition, Radar, Communication networks

\item Introduction to Probability Theory

Experiment, Sample Space, Events

Axioms of Probability

Corollaries and Propositions from axioms

Assigning probability (Classical \& Relative frequency) (mentioned but not shown in the visible pages)

\item Joint Probability (mentioned in outline but not included in these 20 pages)

Motivation, notation, concept

Card deck example

Costume party example

\item Conditional Probability (mentioned in outline but not included in these 20 pages)

Motivation, notation, concept

Cards without replacement

Poker

Missing key
\end{enumerate}

So in this PDF portion, we can fully study:

intro definitions + axioms + corollaries + propositions.

\section*{1) Why should we learn Probability? (Computing Motivation)}

The lecture first explains probability through engineering applications.

\subsection*{1.1 Speech Recognition System (pages 3–4)}

Conceptual idea:

Speech recognition systems store a vocabulary like:

Hello

Yes

No

Bye

And for each word, the system stores templates (waveforms).

Why probability is needed?

Because the input speech can vary due to:

Noise

Different speakers

Interruptions

So the system must decide:

Which word has the highest probability given the observed signal?

That is exactly a probability-based decision.

\subsection*{1.2 Radar System (page 5)}

Radar tries to detect whether an aircraft is present.

The slide shows two hypotheses:

$H_0$ : only noise

$H_1$ : signal + noise

And it shows errors:

False Alarm

Miss Detection

So probability helps measure these risks.

\subsection*{1.3 Communication Networks (page 6)}

Networks have uncertainty like:

arrival of packets

delay (latency)

QoS variations

Probability is used to model and predict these uncertain behaviors.

\section*{2) Probability Basics: Experiment, Outcome, Event, Sample Space}

This is the foundation of the whole course.

\subsection*{2.1 Experiment (E)}

Experiment (E):

A procedure we perform, that produces some result.

Example (given):

Tossing a coin five times → $E_5$

\subsection*{2.2 Outcome ($\xi$)}

Outcome ($\xi$):

A possible result of an experiment.

Example (given):

One outcome of $E_5$:

$\xi_1 = HHTHT$

Beginner clarity:

Outcome = one exact result

Like a single sequence in 5 tosses.

\subsection*{2.3 Event}

Event (Any letter):

A certain set of outcomes of an experiment.

Example (given):

Event $C$ in experiment $E_5$:

$C = \{ \text{all outcomes consisting of an even number of heads} \}$

Intuition:

Event is like a “condition”.

Not one outcome, but a collection.

\subsection*{2.4 Sample Space (S)}

Sample Space (S):

The collection or set of all possible distinct outcomes of an experiment.

The lecture emphasizes sample space outcomes are:

Mutually Exclusive

Collectively Exhaustive

(A) Mutually Exclusive

You can get heads or tails, but not both.

So outcomes do not overlap.

(B) Collectively Exhaustive

You cannot get anything other than heads or tails.

So all possibilities are covered.

\subsection*{2.5 Types of Sample Space (page 5)}

The lecture says sample space can be:

1) Discrete

Finite outcomes.

Example:

$S = \{H,T\}$

2) Countably Infinite

Infinite outcomes but listable.

Example from lecture:

Flip a coin until tails occurs

Possible outcomes:

T

HT

HHT

HHHT

…

3) Continuous

Uncountably many outcomes.

Example from lecture:

Random number generator in $[0,1)$

$S = \{x : 0 \le x < 1\}$

\section*{3) Probability Definition}

Probability:

It is a measure of the likelihood of various events

OR

It is a function of an event that produces a numerical quantity that measures the likelihood of that event.

\section*{4) Axioms of Probability}

The lecture defines axioms as:

self-evident statements requiring no proof.

Axiom 1: Range

For any event $A$:

\[
0 \le Pr(A) \le 1
\]

Axiom 2: Sample space probability

\[
Pr(S) = 1
\]

Axiom 3: Additivity for mutually exclusive events

(a) Two disjoint events

If:

\[
A \cap B = \emptyset
\]

Then:

\[
Pr(A \cup B) = Pr(A) + Pr(B)
\]

(b) Infinite disjoint events

\[
Pr\left(\bigcup_{i=1}^{\infty} A_i\right)
=
\sum_{i=1}^{\infty} Pr(A_i)
\]

\section*{5) Corollary from Axioms}

Corollary 2.1:

For a finite number $M$ of mutually exclusive sets:

\[
Pr\left(\bigcup_{i=1}^{M} A_i\right)
=
\sum_{i=1}^{M} Pr(A_i)
\]

Explanation:

This is basically Axiom 3, but for finite unions.

\section*{6) Propositions from Axioms}

Proposition 2.1: Complement Rule

\[
Pr(A^c) = 1 - Pr(A)
\]

Stepwise derivation (very exam important)

$A$ and $A^c$ cannot happen together:

\[
A \cap A^c = \emptyset
\]

Together they cover the full sample space:

\[
A \cup A^c = S
\]

From Axiom 2:

\[
Pr(S) = 1
\]

From Axiom 3(a):

\[
Pr(A \cup A^c) = Pr(A) + Pr(A^c)
\]

Substitute:

\[
1 = Pr(A) + Pr(A^c)
\]

Rearrange:

\[
Pr(A^c) = 1 - Pr(A)
\]

Proposition 2.2: Monotonicity Rule

If $A \subset B$, then:

\[
Pr(A) \le Pr(B)
\]

\section*{7) Simple Diagrams (for understanding)}

Event inside sample space

\begin{verbatim}
+-------------------+
|         S         |
|   +----------+    |
|   |    A     |    |
|   +----------+    |
+-------------------+
\end{verbatim}

Complement

\begin{verbatim}
+-------------------+
|         S         |
|   +----------+    |
|   |    A     |    |
|   +----------+    |
|  rest = A^c       |
+-------------------+
\end{verbatim}

\section*{8) Practice Examples (Beginner Friendly)}

Example 1: Complement

If:

\[
Pr(A) = 0.65
\]

Then:

\[
Pr(A^c) = 1 - 0.65 = 0.35
\]

Example 2: Additivity (Mutually exclusive)

Die roll:

$A = \{1\}$

$B = \{2\}$

\[
Pr(A)=\frac{1}{6}, \quad Pr(B)=\frac{1}{6}
\]

\[
Pr(A \cup B) = \frac{1}{6} + \frac{1}{6} = \frac{2}{6}
\]

Example 3: Monotonicity

Let:

$A = \{\text{roll 6}\}$

$B = \{\text{roll even}\} = \{2,4,6\}$

Then:

\[
A \subset B \Rightarrow Pr(A) \le Pr(B)
\]

\section*{9) Final Summary (Very Concise)}

This lecture portion establishes:

the building blocks of probability (experiment → outcomes → events → sample space)

the three probability axioms

how to derive key probability results like:

complement rule

monotonicity rule

additivity for finite unions

These are the base tools required for joint and conditional probability later.

\section*{10) Key Formula Sheet (Exam Ready)}

Axioms:

\[
0 \le Pr(A) \le 1
\]

\[
Pr(S)=1
\]

If $A \cap B = \emptyset$:

\[
Pr(A \cup B)=Pr(A)+Pr(B)
\]

Infinite union:

\[
Pr\left(\bigcup_{i=1}^{\infty} A_i\right)
=
\sum_{i=1}^{\infty} Pr(A_i)
\]

Corollary (finite union):

\[
Pr\left(\bigcup_{i=1}^{M} A_i\right)
=
\sum_{i=1}^{M} Pr(A_i)
\]

Complement:

\[
Pr(A^c)=1-Pr(A)
\]

Monotonicity:

If $A \subset B$:

\[
Pr(A) \le Pr(B)
\]

\section*{11) Probable Exam Questions}

Very likely theory questions:

Define experiment, outcome, event, sample space

Explain mutually exclusive and collectively exhaustive

State the 3 axioms of probability

Derive $Pr(A^c)=1-Pr(A)$

Explain why $A \subset B \Rightarrow Pr(A) \le Pr(B)$

Very likely numericals:

Compute complement probabilities

Compute union probability of disjoint events

Identify sample space type (discrete/countably infinite/continuous)

\end{document}
